% !TeX root = ../../main.tex
% ============================================================
% Appendix C: Enemy（新建，仿照 A_Hierarchy / B_Player 的结构）
% 说明：本文件由 main.tex 合并进主文档，且位于 main.tex 的
%       \appendix 之后，因此 \section 会自动编号为 Appendix C。
%       内容忠实来自 策划案.docx 原文（英文）。
%       路径中的反斜杠用 \bs （在 main.tex 中定义为 \textbackslash）。
% ============================================================

\section{Enemy}   % 自动显示为 "Appendix C: Enemy"
\subsection{EnemyMovement Direction Submodule}\label{app:c:enemy-movement}

\par \hspace{1.5em} Relative Path: \texttt{\bs Assets\bs EnemyUnit\bs PublicScript\bs EnemyMovement\bs Direction}

\begin{enumerate}
	\item Direction Data Maintainance: EnemyDirectorCalculator.cs.
	\item Avoid the obstacle: EnemyAvoidance.cs.
	\item Seperate the enemies movement(do not let enemies engaged together): EnemySeperation.cs.
	\item Decide the state whether to ``Patrol'' or ``Chase''.
	\item Use the BFS algorithm to optimize the chasing algorithm. The algorithm placed in the ``PathFingingAlgorithm'' folder and the algorithm's description placed in \ref{app:c:bfs}.
\end{enumerate}

\subsection{PathFinding Algorithm}\label{app:c:bfs}
\par\hspace{1.5em}The BFS(Broad Field Searching) algorithm will increase the enemy chasing logic, preventing enemy chasing like a noob.
	The BFS algorithm achieved by following steps:
\begin{enumerate}
	\item Scan the whole map into states that combined with 'can walk', 'cannot walk'.
	\item Applied the BFS algorithm into the EnemyMovement's Direction data flow to adjust the moving direction so that enemy will follow the correct path to chase the player.
	\item\textbf{The BFS algorithm's arguments have written in the BFS\_Arguments.cs.}
\end{enumerate}